\section{结论与后续方向}\label{sec:conclusions}

在本讲义中，我们首先讨论了 $\cN{=}2$ 超对称规范理论的拉格朗日量，然后研究了含有不同味数的 $\SU(2)$ 规范理论的库仑支和希格斯支。在这一过程中，Seiberg-Witten 曲线和紫外曲线这两个相关概念发挥了非常重要的作用。随后，我们分析了精细调节库仑支真空期望值或恰好边缘耦合参数时会发生什么。有时，这一极限由一个对偶的弱耦合规范理论描述，含有四个味的 $\SU(2)$ 理论便是如此。然而，在大多数情况下，我们看到最终得到的是新的超共形场论，属于 Argyres-Douglas 型或 Gaiotto 型。

例如，我们在第 \ref{sec:rank-1SCFT} 小节中见到了理论 $AD_{N_f=1,2,3}(\SU(2))$ 和 $MN(E_{6,7,8})$，在第 \ref{sec:generalAD} 小节中见到了理论 $X_N$ 和 $Y_N$，在第 \ref{sec:partII} 小节中见到了 $R_N$，在第 \ref{sec:application} 小节中见到了 $T_N$。越来越多的 $\cN{=}2$ 超共形理论正在被发现，例如见文献 \cite{Cecotti:2010fi,Cecotti:2012jx,Cecotti:2013lda}。这意味着，要充分理解 $\cN{=}2$ 超对称系统之间的相互关系，就不能只关注由矢量多重态和超多重态组成的 $\cN{=}2$ 理论。

本讲义所涉及的主题，只是 $\cN{=}2$ 动力学研究这座巨大冰山的一角，还有许多进一步的研究方向。下面列举其中一些。\footnote{作者在这里并未力求穷尽和全面，只引用了少数近期文献。如有人提出请求，作者很乐意补充更多引文。} 首先，可以把一个 $\cN{=}2$ 理论放在非平凡流形上：
% [Tutor Note: p.169 原文 "a nontrival manifold" 中 "nontrival" 应为 "nontrivial"，系拼写笔误；译文按本意作“非平凡流形”。]
\begin{itemize}
\item 利用拓扑扭曲\index{TopologicalTwist@拓扑扭曲|hyperpage}，可以把它放在任意流形上 \cite{Witten:1988ze}。当流形紧致时，配分函数等价于数学家所说的 Donaldson 不变量。将 Seiberg-Witten 解应用于纯 $\SU(2)$ 理论，Witten 引入了一个新的数学不变量，如今称为 Seiberg-Witten 不变量\index{SeibergWittenInvariant@Seiberg-Witten 不变量|hyperpage} \cite{Witten:1994cg}；二十年前，它给四维微分几何带来了一场革命。
\item 可以把它放在 $S^1$ 上。此时理论在有效意义上是三维的。正如文献 \cite{Seiberg:1996nz} 首次分析的那样，作为三维理论的库仑支，自然地是在四维理论库仑支之上的一个纤维化。三维库仑支是超凯勒的，并具有有限自由度经典可积系统的结构。这一可积系统最初在文献 \cite{Donagi:1995cf} 中引入。关于现代发展，例如见文献 \cite{Gaiotto:2008cd}。
\item 可以把它放在所谓的 $\Omega$ 背景上。非常粗略地说，这涉及使 $\bR^4$ 上的整个欧氏系统绕原点作强制旋转。时空在有效意义上是紧致的，因而可以定义配分函数，通常称为 Nekrasov 配分函数。关于近期的全面讨论，例如见文献 \cite{Nekrasov:2012xe}。文献 \cite{Nekrasov:2009rc} 发现，在某种极限情形下，它给出一个量子化可积系统，即 Donagi-Witten 可积系统的量子化版本。
\item 可以把它放在圆的或形变的 $S^4$ 上。时空是紧致的，配分函数可以精确计算，例如见文献 \cite{Pestun:2007rz,Hama:2012bg,Nosaka:2013cpa}。还已知该配分函数与紫外曲线上的二维共形场论有关，例如见文献 \cite{Alday:2009aq,Wyllard:2009hg}。
\item 可以把它放在 $S^1\times S^3$ 上。配分函数称为超共形指标\index{SuperconformalIndex@超共形指标|hyperpage}，并给出紫外曲线上的二维拓扑场论。它还与多种重要的正交多项式有深刻联系，例如见文献 \cite{Gadde:2011uv,Gaiotto:2012xa}。
\item 也可以考虑其他背景。关于 $S^2\times S^1\times\bR$，见文献 \cite{Luo:2013nxa}。文献 \cite{Closset:2013sxa} 研究了 $T^2\times S^2$ 上的 $\cN{=}1$ 理论，当然也可以在那里类似地考虑 $\cN{=}2$ 系统。
\end{itemize}
其次，可以研究这些理论的动力学激发和从外部引入的算符：
\begin{itemize}
\item 我们已经看到，如何从六维构造读出饱和 BPS 界的粒子种类数。这个数是整数，因此通常不能改变，但它确实会在库仑支上的某些位置发生跳跃。这称为跨墙\index{WallCrossing@跨墙|hyperpage}，是一个受到深入研究的领域，例如见文献 \cite{Gaiotto:2009hg}。所得的谱通常可以用一种称为 BPS 箭图\index{BPSQuiver@BPS 箭图|hyperpage}的图来概括。这一观点最初在文献 \cite{Denef:2002ru} 中的 $\cN{=}2$ 超引力背景下引入。关于较新的发展，例如见文献 \cite{Denef:2007vg,Cecotti:2011rv,Manschot:2012rx}。
\item 不引入动力学粒子，也可以引入外部对象的世界线。它们称为线算符\index{LineOperator@线算符|hyperpage}。例如见文献 \cite{Drukker:2009tz,Gaiotto:2010be}。
\item 一旦允许引入外部线算符，就没有理由不引入更高维的外部对象。当它们具有两个时空维度时，称为面算符\index{SurfaceOperator@面算符|hyperpage}。Seiberg-Witten 曲线可以内禀地定义为一个面算符的红外模空间 \cite{Gaiotto:2009fs}。另一篇值得研读的近期有趣论文是文献 \cite{Gaiotto:2013sma}。
\item 进而可以考虑具有三个时空维度的对象。这就是外部畴壁\index{DomainWall@畴壁|hyperpage}。关于近期的研究，例如见文献 \cite{Dimofte:2013lba}。
\end{itemize}
关于这些主题，综述 \cite{MooreReview} 是极好的资料来源，尽管该综述本身是面向数学家的。

第三，本讲义所描述的方法还不足以求解任意的 $\cN{=}2$ 规范理论。许多四维 $\cN{=}2$ 理论确实来自六维 $\cN{=}(2,0)$ 理论，但目前也有许多理论并非如此。因此，还应该研究其他方法。
\begin{itemize}
\item 六维构造本身还需要进一步发展。关于驯服穿孔的进一步讨论，例如见文献 \cite{Chacaltana:2010ks,Chacaltana:2012ch,Chacaltana:2013oka}；关于野性穿孔，更多内容见文献 \cite{Xie:2012hs,Xie:2013jc}。
\item 四维 $\cN{=}2$ 理论的六维构造，总能提升为非紧致 Calabi-Yau 流形上的 IIB 型弦论，该流形是紫外曲线之上的一个纤维化。即便这个非紧致 Calabi-Yau 流形不是曲线之上的纤维化，其上的 IIB 型弦论也常常实现一个四维 $\cN{=}2$ 场论，这为求解 $\cN{=}2$ 系统提供了另一种途径，例如见文献 \cite{Tachikawa:2011yr,Cecotti:2012sf}。
\item 对于一般规范群 $G$，$\cN{=}2^*$ 理论，即通过伴随超多重态的质量项使 $\cN{=}4$ 超杨--米尔斯理论发生形变而得到的理论，早已得到求解 \cite{D'Hoker:1999ft}。有些令人惊讶的是，当 $G\neq\SU(N)$ 时，目前还不知道这些解的明确弦论或 M 理论构造。这清楚地表明，我们目前的理解还多么初步。
\end{itemize}

第四，$\cN{=}2$ 理论有许多性质为所有已知例子所满足，但目前还没有办法推导它们。设计研究这些性质的新方法，应当会富有成效。下面列出这一方向上的几个问题。
\begin{itemize}
\item $\cN{=}2$ 规范理论库仑支上的手征算符显然总是自由生成的。例如，在 $\SU(N)$ 规范理论中，它们由 $\tr\phi^k$（$k=2,\ldots,N$）生成，这些算符之间没有非平凡关系。从已有例子来看，由六维构造得到的所有非拉格朗日理论仍满足这一性质：库仑支算符是自由生成的。作者猜想，这实际上是一条适用于每个 $\cN{=}2$ 超对称系统的定理。
\item 文献 \cite{Buican:2013ica} 论证了，在两个 $\cN{=}2$ 超共形场论之间的重整化群流中，中心荷 $a$ 的变化存在一个非零下界。能否更严格地推导这一事实？
% [Tutor Note: p.171 原文 "Is there are more rigorous derivation of this fact?" 中 "there are" 应为 "there a"（即 "Is there a more rigorous derivation …"），系笔误；译文按本意译出。]
\item 能否刻画整个 $\cN{=}2$ 理论谱系本身？作一个类比，考虑 $\SU(2)$ 的所有表示。如果只允许直和，就需要所有不可约表示，才能构造所有可能的表示。如果还允许张量积以及取出一个不可约直和分量，那么只需要二维不可约表示，就能生成所有其他表示。

对于 $\cN{=}2$ 理论，也可以提出类似的问题。如果只允许弱规范化，需要哪些广义物质内容，即超多重态和其他“不可约”的强耦合理论，才能生成所有 $\cN{=}2$ 理论？如果还允许强耦合极限、S 对偶以及分解为组成部分，那么需要多少？通过六维、弦论或 M 理论构造，能够得到这些理论的多大“百分比”？

在某种技术意义下完备的 $\cN{=}2$ 理论，已在文献 \cite{Cecotti:2011rv} 中得到分类；$\cN{=}2$ 弱耦合规范理论，已在文献 \cite{Bhardwaj:2013qia} 中得到分类。不过，这只是进入所有可能的 $\cN{=}2$ 理论所构成的广阔空间的两个小小步骤。本讲义中见到的理论汇总于附录 \ref{sec:zoo}。
\end{itemize}

最后，作者想强调，即便是下列看起来如此简单的规范理论：
\begin{itemize}
\item 带有一个三指标反对称张量表示超多重态的 $\cN{=}2$ 超对称 $\SU(7)$ 规范理论，或
\item 带有一个有质量的三基本表示完整超多重态 $(Q_{aiu},\tilde Q^{aiu})$ 的 $\cN{=}2$ 超对称 $\SU(2)^3$ 规范理论，
\end{itemize}
也还没有得到求解。对于第一个求出这两个理论中任意一个的解的人，作者很乐意在柏校区的寿司店请他吃一顿晚餐。如文献 \cite{Bhardwaj:2013qia} 所列，还有许多其他 $\cN{=}2$ 规范理论没有已知的解。因此，这个领域仍应被看作大有可为。

希望读到这里的读者，至少已经具备了一定的基础，可以着手研读这些以及其他关于 $\cN{=}2$ 超对称理论的近期论文。如果他们继续学习，并为拓展研究前沿作出贡献，作者将感到十分欣慰。
